% Generated draft: requires mathematical and source review.
\documentclass[11pt]{article}
\usepackage[margin=25mm]{geometry}
\usepackage{fontspec}
\setmainfont{Latin Modern Roman}
\usepackage{amsmath,amssymb,mathtools,unicode-math}
\setmathfont{Latin Modern Math}
\usepackage{enumitem,needspace,xurl,hyperref}
\setlength{\parindent}{0pt}
\setlength{\parskip}{5pt}
\setlength{\emergencystretch}{3em}
\newcommand{\N}{\mathbb{N}}
\newcommand{\Z}{\mathbb{Z}}
\newcommand{\Q}{\mathbb{Q}}
\newcommand{\R}{\mathbb{R}}
\newcommand{\C}{\mathbb{C}}
\newcommand{\F}{\mathbb{F}}
\newcommand{\A}{\mathbb{A}}
\newcommand{\PP}{\mathbb{P}}
\newcommand{\E}{\mathbb{E}}
\newcommand{\eps}{\varepsilon}
\newcommand{\dd}{\,\mathrm{d}}
\begin{document}
\section*{UnsolvedMath ID: 10600054}
\textbf{Source problem code:} catalogue entry
\subsection{Definitions and mathematical statement}
Let \(\mathcal G\) be the set of graph-link equivalence classes in the source's graph-Reidemeister theory, and let \(\mathcal R\subseteq\mathcal G\) be the realizable classes. In the customary chord-diagram interpretation, a realizable representative is an appropriately labelled graph obtained as an intersection graph of a chord diagram: vertices correspond to chords and adjacency means that their endpoints alternate. The excerpt omits the labels, graph moves and exact realizability convention; in particular it must be stated whether membership of \(\mathcal R\) means that a class has at least one realizable representative.

The source asks to construct a projection onto realizable graph-links. The standard retraction reading is a map
\[P:\mathcal G\to\mathcal R,\qquad P(r)=r\quad(r\in\mathcal R).\]
With inclusion \(i:\mathcal R\hookrightarrow\mathcal G\), this says \(P\circ i=\operatorname{id}_{\mathcal R}\), and \(i\circ P\) is idempotent. A rule on graph representatives must respect all graph-link equivalences to define such a map on classes.

The intended additional meaning of ``construct'' or ``projection'' is missing. A nonempty subset of a set admits a set-theoretic retraction with no useful geometric content, so the substantive problem needs, for example, an effective diagram rule, compatibility with graph operations, a prescribed relation to the original graph, or some other stated property. None of these examples is imposed as a requirement not present in the source. Existence of a set map, decidability of realizability of classes, and a local/natural projection are different tasks; the excerpt does not choose among them.
\subsection{Short English statement}
Construct a well-defined projection from graph-link classes onto the realizable ones, fixing realizable classes. The problem needs a further constructive or structural condition to distinguish its intended projection from a bare set retraction.
\subsection{Sources}
\begin{itemize}
\item Local UnsolvedMath registry, permanent ID \(10600054\), exact source code \texttt{catalogue entry}.
\item \emph{Unsolved Problems in Virtual Knot Theory and Combinatorial Knot Theory}, revised survey, problem-list item 54: \url{https://arxiv.org/abs/1409.2823}. Registry access 2026-07-29; source-tex extraction. No separate primary construction reference is supplied.
\end{itemize}
\end{document}
